% testing_EN.tex — WAC (Windows Artefacts Collector) testing guide, English edition.
% The French edition is testing_FR.tex, in this folder: both are maintained together,
% section for section.
%
% Build (LuaLaTeX, as the editor's recipe; build artefacts removed after):
%     ../build-docs.sh testing_EN.tex
% The cover page is ../cover_tests_EN.png, derived from ../cover_EN.png by
% ../retitle-cover.py. The layout is ../wac-style.tex, shared with the user manual.

\documentclass[11pt,a4paper,oneside]{report}

% ---------------------------------------------------------------- language and labels
\usepackage[british]{babel}
\newcommand{\wacdoctitle}{Testing guide}
\newcommand{\wacpageof}{of}
\newcommand{\wacnotetitle}{Keep in mind}
\newcommand{\wacwarningtitle}{Warning}
% Layout shared with every WAC document.
% wac-style.tex — layout shared by every WAC document (user manuals, testing
% guides), French and English editions alike: one copy, so that they cannot
% drift apart.
%
% A document sets, BEFORE % wac-style.tex — layout shared by every WAC document (user manuals, testing
% guides), French and English editions alike: one copy, so that they cannot
% drift apart.
%
% A document sets, BEFORE % wac-style.tex — layout shared by every WAC document (user manuals, testing
% guides), French and English editions alike: one copy, so that they cannot
% drift apart.
%
% A document sets, BEFORE \input{../wac-style.tex}:
%   \usepackage[<language>]{babel}
%   \newcommand{\wacdoctitle}{...}      running header: "WAC — <title>"
%   \newcommand{\wacpageof}{...}        footer: "Page N <of> M"
%   \newcommand{\wacnotetitle}{...}     default title of the "note" box
%   \newcommand{\wacwarningtitle}{...}  default title of the "attention" box
% and AFTER it its \hypersetup (PDF title, author, subject).

% ---------------------------------------------------------------- fonts, language
\usepackage{fontspec}
\setmainfont{Noto Sans}
\setmonofont{DejaVu Sans Mono}[Scale=0.84]

% ---------------------------------------------------------------- layout
\usepackage[top=26mm,bottom=26mm,left=24mm,right=22mm,headheight=14pt]{geometry}
\usepackage{graphicx}
\graphicspath{{../}}
\usepackage{xcolor}
\usepackage{tikz}
\usepackage{booktabs,tabularx,longtable,array}
\usepackage{enumitem}
\usepackage{fancyhdr}
\usepackage{titlesec}
\usepackage[most]{tcolorbox}
\usepackage{listings}
\usepackage{url}
\usepackage{lastpage}
\usepackage[hidelinks,unicode,pdfpagelabels]{hyperref}

% Colours sampled from the cover page
\definecolor{marine}{HTML}{0A1A33}
\definecolor{bleu}{HTML}{2F6DB5}
\definecolor{accent}{HTML}{1F66D1}
\definecolor{fond}{HTML}{EEF3FA}
\definecolor{alerte}{HTML}{B3261E}
\definecolor{fondalerte}{HTML}{FBEDEC}
\definecolor{gris}{HTML}{5B6573}

\setlength{\parindent}{0pt}
\setlength{\parskip}{0.55em}
\linespread{1.08}
\setlist{itemsep=0.25em,topsep=0.35em}
\urlstyle{tt}
\renewcommand{\arraystretch}{1.25}
% Left-aligned table columns: justified narrow columns opened wide gaps
% between words.
\newcolumntype{L}[1]{>{\raggedright\arraybackslash}p{#1}}
% Monospaced file and field names cannot be hyphenated: let the line's spaces
% stretch instead.
\setlength{\emergencystretch}{3em}

% Headings
\titleformat{\chapter}[display]
  {\color{marine}\bfseries\raggedright}
  {\color{bleu}\large\MakeUppercase{\chaptertitlename}~\thechapter}
  {0.3em}
  {\Huge}[{\color{accent}\titlerule[2pt]}]
\titlespacing*{\chapter}{0pt}{0pt}{1.6em}
\titleformat{\section}{\color{marine}\Large\bfseries}{\color{bleu}\thesection}{0.8em}{}
\titleformat{\subsection}{\color{marine}\large\bfseries}{\color{bleu}\thesubsection}{0.7em}{}
\titleformat{\subsubsection}{\color{bleu}\normalsize\bfseries}{}{0em}{}

% Headers and footers
\pagestyle{fancy}
\fancyhf{}
\fancyhead[L]{\small\color{gris}WAC — \wacdoctitle}
\fancyhead[R]{\small\color{gris}\leftmark}
% Footer: "Page N of M" (the word set by \wacpageof), readable. A bare, small, light-grey number went
% unnoticed.
\newcommand{\piedpage}{\small\color{marine}\textbf{Page \thepage}\ \wacpageof\ \pageref*{LastPage}}
\fancyfoot[C]{\piedpage}
\renewcommand{\footrulewidth}{0.4pt}
\renewcommand{\footrule}{\hbox to\headwidth{\color{fond}\leaders\hrule height \footrulewidth\hfill}}
\renewcommand{\headrulewidth}{0.4pt}
\renewcommand{\headrule}{\hbox to\headwidth{\color{fond}\leaders\hrule height \headrulewidth\hfill}}
\renewcommand{\chaptermark}[1]{\markboth{#1}{}}
\fancypagestyle{plain}{\fancyhf{}\fancyfoot[C]{\piedpage}\renewcommand{\headrulewidth}{0pt}}

% Boxes
\newtcolorbox{note}[1][\wacnotetitle]{enhanced,breakable,colback=fond,colframe=bleu,
  boxrule=0pt,leftrule=3pt,arc=0pt,left=8pt,right=8pt,top=6pt,bottom=6pt,
  title={\bfseries #1},coltitle=bleu,colbacktitle=fond,attach title to upper={\par\smallskip}}
\newtcolorbox{attention}[1][\wacwarningtitle]{enhanced,breakable,colback=fondalerte,colframe=alerte,
  boxrule=0pt,leftrule=3pt,arc=0pt,left=8pt,right=8pt,top=6pt,bottom=6pt,
  title={\bfseries #1},coltitle=alerte,colbacktitle=fondalerte,attach title to upper={\par\smallskip}}

% Code
\lstset{basicstyle=\ttfamily\small,columns=fullflexible,keepspaces=true,
  backgroundcolor=\color{fond},frame=l,framerule=3pt,rulecolor=\color{bleu},
  xleftmargin=9pt,framexleftmargin=6pt,aboveskip=0.8em,belowskip=0.8em,
  breaklines=true,breakatwhitespace=false,showstringspaces=false,
  literate={é}{{\'e}}1 {è}{{\`e}}1 {à}{{\`a}}1 {ê}{{\^e}}1 {ô}{{\^o}}1 {ç}{{\c c}}1
           {É}{{\'E}}1 {—}{{---}}1 {’}{{'}}1 {…}{{...}}1 {«}{{\guillemotleft}}1 {»}{{\guillemotright}}1}

% Breakable monospace. A path or a field name in \texttt is one word that no
% hyphenation pattern cuts: in a table cell it ran 190 pt past the column
% (ProgramData\Microsoft\NetFramework\BreadcrumbStore). As the \seqsplit of
% the EventsManager theme's tables, a line may break inside it, without a
% hyphen that would read as part of the name: after \ / . _ - : and between a
% lowercase letter or digit and a capital (AmcacheSha1|Covers|Bytes). Only
% inside \texttt — the classes are set once, the breaks switched on there.
\newXeTeXintercharclass\wacBreakAfter
\newXeTeXintercharclass\wacUpper
\newXeTeXintercharclass\wacLower
\XeTeXcharclass`\\=\wacBreakAfter
\XeTeXcharclass`\/=\wacBreakAfter
\XeTeXcharclass`\.=\wacBreakAfter
\XeTeXcharclass`\_=\wacBreakAfter
\XeTeXcharclass`\-=\wacBreakAfter
\XeTeXcharclass`\:=\wacBreakAfter
\count255=`A \loop\XeTeXcharclass\count255=\wacUpper\ifnum\count255<`Z\advance\count255 by 1\repeat
\count255=`a \loop\XeTeXcharclass\count255=\wacLower\ifnum\count255<`z\advance\count255 by 1\repeat
\count255=`0 \loop\XeTeXcharclass\count255=\wacLower\ifnum\count255<`9\advance\count255 by 1\repeat
\XeTeXinterchartoks\wacBreakAfter\wacUpper={\allowbreak}
\XeTeXinterchartoks\wacBreakAfter\wacLower={\allowbreak}
\XeTeXinterchartoks\wacBreakAfter 0={\allowbreak}
\XeTeXinterchartoks\wacLower\wacUpper={\allowbreak}
% \textbackslash is no character token: it carries its own break.
\newcommand{\wacbackslash}{\char92\allowbreak}
\let\wacplaintexttt\texttt
\renewcommand{\texttt}[1]{\wacplaintexttt{\XeTeXinterchartokenstate=1\let\textbackslash\wacbackslash#1}}

% Shortcuts
\newcommand{\wac}{\textsc{wac}}
\newcommand{\opt}[1]{\texttt{-{}-#1}}
\newcommand{\champ}[1]{\texttt{#1}}
\newcommand{\fichier}[1]{\texttt{#1}}

% ============================================================ COVER PAGE
% The cover images are 2:3 (1024 x 1536), taller than an A4 page once widened
% to it: centred, they lost 9 mm at the top — the CEC logo — and 9 mm at the
% bottom. Raised 4 mm above the top of the page instead: the logo, which starts
% 6 mm below the image's edge, stays whole, and the bottom row of text keeps
% 7 mm to the page's edge; only a decorative corner below it is cut.
% Usage: \coverpage{cover_FR.png}
\newcommand{\coverpage}[1]{%
  \begin{titlepage}
  \thispagestyle{empty}
  \begin{tikzpicture}[remember picture,overlay]
    \node[inner sep=0pt,anchor=north,yshift=4mm] at (current page.north)
      {\includegraphics[width=\paperwidth]{#1}};
  \end{tikzpicture}
  \end{titlepage}}
:
%   \usepackage[<language>]{babel}
%   \newcommand{\wacdoctitle}{...}      running header: "WAC — <title>"
%   \newcommand{\wacpageof}{...}        footer: "Page N <of> M"
%   \newcommand{\wacnotetitle}{...}     default title of the "note" box
%   \newcommand{\wacwarningtitle}{...}  default title of the "attention" box
% and AFTER it its \hypersetup (PDF title, author, subject).

% ---------------------------------------------------------------- fonts, language
\usepackage{fontspec}
\setmainfont{Noto Sans}
\setmonofont{DejaVu Sans Mono}[Scale=0.84]

% ---------------------------------------------------------------- layout
\usepackage[top=26mm,bottom=26mm,left=24mm,right=22mm,headheight=14pt]{geometry}
\usepackage{graphicx}
\graphicspath{{../}}
\usepackage{xcolor}
\usepackage{tikz}
\usepackage{booktabs,tabularx,longtable,array}
\usepackage{enumitem}
\usepackage{fancyhdr}
\usepackage{titlesec}
\usepackage[most]{tcolorbox}
\usepackage{listings}
\usepackage{url}
\usepackage{lastpage}
\usepackage[hidelinks,unicode,pdfpagelabels]{hyperref}

% Colours sampled from the cover page
\definecolor{marine}{HTML}{0A1A33}
\definecolor{bleu}{HTML}{2F6DB5}
\definecolor{accent}{HTML}{1F66D1}
\definecolor{fond}{HTML}{EEF3FA}
\definecolor{alerte}{HTML}{B3261E}
\definecolor{fondalerte}{HTML}{FBEDEC}
\definecolor{gris}{HTML}{5B6573}

\setlength{\parindent}{0pt}
\setlength{\parskip}{0.55em}
\linespread{1.08}
\setlist{itemsep=0.25em,topsep=0.35em}
\urlstyle{tt}
\renewcommand{\arraystretch}{1.25}
% Left-aligned table columns: justified narrow columns opened wide gaps
% between words.
\newcolumntype{L}[1]{>{\raggedright\arraybackslash}p{#1}}
% Monospaced file and field names cannot be hyphenated: let the line's spaces
% stretch instead.
\setlength{\emergencystretch}{3em}

% Headings
\titleformat{\chapter}[display]
  {\color{marine}\bfseries\raggedright}
  {\color{bleu}\large\MakeUppercase{\chaptertitlename}~\thechapter}
  {0.3em}
  {\Huge}[{\color{accent}\titlerule[2pt]}]
\titlespacing*{\chapter}{0pt}{0pt}{1.6em}
\titleformat{\section}{\color{marine}\Large\bfseries}{\color{bleu}\thesection}{0.8em}{}
\titleformat{\subsection}{\color{marine}\large\bfseries}{\color{bleu}\thesubsection}{0.7em}{}
\titleformat{\subsubsection}{\color{bleu}\normalsize\bfseries}{}{0em}{}

% Headers and footers
\pagestyle{fancy}
\fancyhf{}
\fancyhead[L]{\small\color{gris}WAC — \wacdoctitle}
\fancyhead[R]{\small\color{gris}\leftmark}
% Footer: "Page N of M" (the word set by \wacpageof), readable. A bare, small, light-grey number went
% unnoticed.
\newcommand{\piedpage}{\small\color{marine}\textbf{Page \thepage}\ \wacpageof\ \pageref*{LastPage}}
\fancyfoot[C]{\piedpage}
\renewcommand{\footrulewidth}{0.4pt}
\renewcommand{\footrule}{\hbox to\headwidth{\color{fond}\leaders\hrule height \footrulewidth\hfill}}
\renewcommand{\headrulewidth}{0.4pt}
\renewcommand{\headrule}{\hbox to\headwidth{\color{fond}\leaders\hrule height \headrulewidth\hfill}}
\renewcommand{\chaptermark}[1]{\markboth{#1}{}}
\fancypagestyle{plain}{\fancyhf{}\fancyfoot[C]{\piedpage}\renewcommand{\headrulewidth}{0pt}}

% Boxes
\newtcolorbox{note}[1][\wacnotetitle]{enhanced,breakable,colback=fond,colframe=bleu,
  boxrule=0pt,leftrule=3pt,arc=0pt,left=8pt,right=8pt,top=6pt,bottom=6pt,
  title={\bfseries #1},coltitle=bleu,colbacktitle=fond,attach title to upper={\par\smallskip}}
\newtcolorbox{attention}[1][\wacwarningtitle]{enhanced,breakable,colback=fondalerte,colframe=alerte,
  boxrule=0pt,leftrule=3pt,arc=0pt,left=8pt,right=8pt,top=6pt,bottom=6pt,
  title={\bfseries #1},coltitle=alerte,colbacktitle=fondalerte,attach title to upper={\par\smallskip}}

% Code
\lstset{basicstyle=\ttfamily\small,columns=fullflexible,keepspaces=true,
  backgroundcolor=\color{fond},frame=l,framerule=3pt,rulecolor=\color{bleu},
  xleftmargin=9pt,framexleftmargin=6pt,aboveskip=0.8em,belowskip=0.8em,
  breaklines=true,breakatwhitespace=false,showstringspaces=false,
  literate={é}{{\'e}}1 {è}{{\`e}}1 {à}{{\`a}}1 {ê}{{\^e}}1 {ô}{{\^o}}1 {ç}{{\c c}}1
           {É}{{\'E}}1 {—}{{---}}1 {’}{{'}}1 {…}{{...}}1 {«}{{\guillemotleft}}1 {»}{{\guillemotright}}1}

% Breakable monospace. A path or a field name in \texttt is one word that no
% hyphenation pattern cuts: in a table cell it ran 190 pt past the column
% (ProgramData\Microsoft\NetFramework\BreadcrumbStore). As the \seqsplit of
% the EventsManager theme's tables, a line may break inside it, without a
% hyphen that would read as part of the name: after \ / . _ - : and between a
% lowercase letter or digit and a capital (AmcacheSha1|Covers|Bytes). Only
% inside \texttt — the classes are set once, the breaks switched on there.
\newXeTeXintercharclass\wacBreakAfter
\newXeTeXintercharclass\wacUpper
\newXeTeXintercharclass\wacLower
\XeTeXcharclass`\\=\wacBreakAfter
\XeTeXcharclass`\/=\wacBreakAfter
\XeTeXcharclass`\.=\wacBreakAfter
\XeTeXcharclass`\_=\wacBreakAfter
\XeTeXcharclass`\-=\wacBreakAfter
\XeTeXcharclass`\:=\wacBreakAfter
\count255=`A \loop\XeTeXcharclass\count255=\wacUpper\ifnum\count255<`Z\advance\count255 by 1\repeat
\count255=`a \loop\XeTeXcharclass\count255=\wacLower\ifnum\count255<`z\advance\count255 by 1\repeat
\count255=`0 \loop\XeTeXcharclass\count255=\wacLower\ifnum\count255<`9\advance\count255 by 1\repeat
\XeTeXinterchartoks\wacBreakAfter\wacUpper={\allowbreak}
\XeTeXinterchartoks\wacBreakAfter\wacLower={\allowbreak}
\XeTeXinterchartoks\wacBreakAfter 0={\allowbreak}
\XeTeXinterchartoks\wacLower\wacUpper={\allowbreak}
% \textbackslash is no character token: it carries its own break.
\newcommand{\wacbackslash}{\char92\allowbreak}
\let\wacplaintexttt\texttt
\renewcommand{\texttt}[1]{\wacplaintexttt{\XeTeXinterchartokenstate=1\let\textbackslash\wacbackslash#1}}

% Shortcuts
\newcommand{\wac}{\textsc{wac}}
\newcommand{\opt}[1]{\texttt{-{}-#1}}
\newcommand{\champ}[1]{\texttt{#1}}
\newcommand{\fichier}[1]{\texttt{#1}}

% ============================================================ COVER PAGE
% The cover images are 2:3 (1024 x 1536), taller than an A4 page once widened
% to it: centred, they lost 9 mm at the top — the CEC logo — and 9 mm at the
% bottom. Raised 4 mm above the top of the page instead: the logo, which starts
% 6 mm below the image's edge, stays whole, and the bottom row of text keeps
% 7 mm to the page's edge; only a decorative corner below it is cut.
% Usage: \coverpage{cover_FR.png}
\newcommand{\coverpage}[1]{%
  \begin{titlepage}
  \thispagestyle{empty}
  \begin{tikzpicture}[remember picture,overlay]
    \node[inner sep=0pt,anchor=north,yshift=4mm] at (current page.north)
      {\includegraphics[width=\paperwidth]{#1}};
  \end{tikzpicture}
  \end{titlepage}}
:
%   \usepackage[<language>]{babel}
%   \newcommand{\wacdoctitle}{...}      running header: "WAC — <title>"
%   \newcommand{\wacpageof}{...}        footer: "Page N <of> M"
%   \newcommand{\wacnotetitle}{...}     default title of the "note" box
%   \newcommand{\wacwarningtitle}{...}  default title of the "attention" box
% and AFTER it its \hypersetup (PDF title, author, subject).

% ---------------------------------------------------------------- fonts, language
\usepackage{fontspec}
\setmainfont{Noto Sans}
\setmonofont{DejaVu Sans Mono}[Scale=0.84]

% ---------------------------------------------------------------- layout
\usepackage[top=26mm,bottom=26mm,left=24mm,right=22mm,headheight=14pt]{geometry}
\usepackage{graphicx}
\graphicspath{{../}}
\usepackage{xcolor}
\usepackage{tikz}
\usepackage{booktabs,tabularx,longtable,array}
\usepackage{enumitem}
\usepackage{fancyhdr}
\usepackage{titlesec}
\usepackage[most]{tcolorbox}
\usepackage{listings}
\usepackage{url}
\usepackage{lastpage}
\usepackage[hidelinks,unicode,pdfpagelabels]{hyperref}

% Colours sampled from the cover page
\definecolor{marine}{HTML}{0A1A33}
\definecolor{bleu}{HTML}{2F6DB5}
\definecolor{accent}{HTML}{1F66D1}
\definecolor{fond}{HTML}{EEF3FA}
\definecolor{alerte}{HTML}{B3261E}
\definecolor{fondalerte}{HTML}{FBEDEC}
\definecolor{gris}{HTML}{5B6573}

\setlength{\parindent}{0pt}
\setlength{\parskip}{0.55em}
\linespread{1.08}
\setlist{itemsep=0.25em,topsep=0.35em}
\urlstyle{tt}
\renewcommand{\arraystretch}{1.25}
% Left-aligned table columns: justified narrow columns opened wide gaps
% between words.
\newcolumntype{L}[1]{>{\raggedright\arraybackslash}p{#1}}
% Monospaced file and field names cannot be hyphenated: let the line's spaces
% stretch instead.
\setlength{\emergencystretch}{3em}

% Headings
\titleformat{\chapter}[display]
  {\color{marine}\bfseries\raggedright}
  {\color{bleu}\large\MakeUppercase{\chaptertitlename}~\thechapter}
  {0.3em}
  {\Huge}[{\color{accent}\titlerule[2pt]}]
\titlespacing*{\chapter}{0pt}{0pt}{1.6em}
\titleformat{\section}{\color{marine}\Large\bfseries}{\color{bleu}\thesection}{0.8em}{}
\titleformat{\subsection}{\color{marine}\large\bfseries}{\color{bleu}\thesubsection}{0.7em}{}
\titleformat{\subsubsection}{\color{bleu}\normalsize\bfseries}{}{0em}{}

% Headers and footers
\pagestyle{fancy}
\fancyhf{}
\fancyhead[L]{\small\color{gris}WAC — \wacdoctitle}
\fancyhead[R]{\small\color{gris}\leftmark}
% Footer: "Page N of M" (the word set by \wacpageof), readable. A bare, small, light-grey number went
% unnoticed.
\newcommand{\piedpage}{\small\color{marine}\textbf{Page \thepage}\ \wacpageof\ \pageref*{LastPage}}
\fancyfoot[C]{\piedpage}
\renewcommand{\footrulewidth}{0.4pt}
\renewcommand{\footrule}{\hbox to\headwidth{\color{fond}\leaders\hrule height \footrulewidth\hfill}}
\renewcommand{\headrulewidth}{0.4pt}
\renewcommand{\headrule}{\hbox to\headwidth{\color{fond}\leaders\hrule height \headrulewidth\hfill}}
\renewcommand{\chaptermark}[1]{\markboth{#1}{}}
\fancypagestyle{plain}{\fancyhf{}\fancyfoot[C]{\piedpage}\renewcommand{\headrulewidth}{0pt}}

% Boxes
\newtcolorbox{note}[1][\wacnotetitle]{enhanced,breakable,colback=fond,colframe=bleu,
  boxrule=0pt,leftrule=3pt,arc=0pt,left=8pt,right=8pt,top=6pt,bottom=6pt,
  title={\bfseries #1},coltitle=bleu,colbacktitle=fond,attach title to upper={\par\smallskip}}
\newtcolorbox{attention}[1][\wacwarningtitle]{enhanced,breakable,colback=fondalerte,colframe=alerte,
  boxrule=0pt,leftrule=3pt,arc=0pt,left=8pt,right=8pt,top=6pt,bottom=6pt,
  title={\bfseries #1},coltitle=alerte,colbacktitle=fondalerte,attach title to upper={\par\smallskip}}

% Code
\lstset{basicstyle=\ttfamily\small,columns=fullflexible,keepspaces=true,
  backgroundcolor=\color{fond},frame=l,framerule=3pt,rulecolor=\color{bleu},
  xleftmargin=9pt,framexleftmargin=6pt,aboveskip=0.8em,belowskip=0.8em,
  breaklines=true,breakatwhitespace=false,showstringspaces=false,
  literate={é}{{\'e}}1 {è}{{\`e}}1 {à}{{\`a}}1 {ê}{{\^e}}1 {ô}{{\^o}}1 {ç}{{\c c}}1
           {É}{{\'E}}1 {—}{{---}}1 {’}{{'}}1 {…}{{...}}1 {«}{{\guillemotleft}}1 {»}{{\guillemotright}}1}

% Breakable monospace. A path or a field name in \texttt is one word that no
% hyphenation pattern cuts: in a table cell it ran 190 pt past the column
% (ProgramData\Microsoft\NetFramework\BreadcrumbStore). As the \seqsplit of
% the EventsManager theme's tables, a line may break inside it, without a
% hyphen that would read as part of the name: after \ / . _ - : and between a
% lowercase letter or digit and a capital (AmcacheSha1|Covers|Bytes). Only
% inside \texttt — the classes are set once, the breaks switched on there.
\newXeTeXintercharclass\wacBreakAfter
\newXeTeXintercharclass\wacUpper
\newXeTeXintercharclass\wacLower
\XeTeXcharclass`\\=\wacBreakAfter
\XeTeXcharclass`\/=\wacBreakAfter
\XeTeXcharclass`\.=\wacBreakAfter
\XeTeXcharclass`\_=\wacBreakAfter
\XeTeXcharclass`\-=\wacBreakAfter
\XeTeXcharclass`\:=\wacBreakAfter
\count255=`A \loop\XeTeXcharclass\count255=\wacUpper\ifnum\count255<`Z\advance\count255 by 1\repeat
\count255=`a \loop\XeTeXcharclass\count255=\wacLower\ifnum\count255<`z\advance\count255 by 1\repeat
\count255=`0 \loop\XeTeXcharclass\count255=\wacLower\ifnum\count255<`9\advance\count255 by 1\repeat
\XeTeXinterchartoks\wacBreakAfter\wacUpper={\allowbreak}
\XeTeXinterchartoks\wacBreakAfter\wacLower={\allowbreak}
\XeTeXinterchartoks\wacBreakAfter 0={\allowbreak}
\XeTeXinterchartoks\wacLower\wacUpper={\allowbreak}
% \textbackslash is no character token: it carries its own break.
\newcommand{\wacbackslash}{\char92\allowbreak}
\let\wacplaintexttt\texttt
\renewcommand{\texttt}[1]{\wacplaintexttt{\XeTeXinterchartokenstate=1\let\textbackslash\wacbackslash#1}}

% Shortcuts
\newcommand{\wac}{\textsc{wac}}
\newcommand{\opt}[1]{\texttt{-{}-#1}}
\newcommand{\champ}[1]{\texttt{#1}}
\newcommand{\fichier}[1]{\texttt{#1}}

% ============================================================ COVER PAGE
% The cover images are 2:3 (1024 x 1536), taller than an A4 page once widened
% to it: centred, they lost 9 mm at the top — the CEC logo — and 9 mm at the
% bottom. Raised 4 mm above the top of the page instead: the logo, which starts
% 6 mm below the image's edge, stays whole, and the bottom row of text keeps
% 7 mm to the page's edge; only a decorative corner below it is cut.
% Usage: \coverpage{cover_FR.png}
\newcommand{\coverpage}[1]{%
  \begin{titlepage}
  \thispagestyle{empty}
  \begin{tikzpicture}[remember picture,overlay]
    \node[inner sep=0pt,anchor=north,yshift=4mm] at (current page.north)
      {\includegraphics[width=\paperwidth]{#1}};
  \end{tikzpicture}
  \end{titlepage}}


\hypersetup{pdftitle={WAC — Testing guide},
  pdfauthor={Aerospace Cyber Defence Centre of Excellence (CEC)},
  pdfsubject={Running the WAC tests by hand and reading their results}}

\begin{document}

% ============================================================ COVER PAGE
\pagenumbering{roman}
\coverpage{cover_tests_EN.png}

% ============================================================ INFORMATION PAGE
\thispagestyle{empty}
\vspace*{\fill}
{\color{marine}\Large\bfseries WAC — Windows Artefacts Collector}\par
{\color{bleu}Testing guide}\par
\medskip
\begin{tabular}{@{}ll@{}}
Software version & 3.0.0 \\
Document edition & October 2026 \\
Publisher & Aerospace Cyber Defence Centre of Excellence (CEC) \\
          & French Air and Space Force Academy, Air Base 701, Salon-de-Provence \\
\end{tabular}

\bigskip
{\small\color{gris}This document explains how to run each of the \wac{} tests
by hand, what each one proves and does not prove, and how to read a failure. It
is meant for whoever changes the code or wants to check a version before using
it. The durations and volumes quoted were measured on the test virtual machine
and on a real collection; they are orders of magnitude. The results of these
tests, and what they guarantee, are in the \emph{Validation report}.}
\vspace*{2cm}
\clearpage

\pagenumbering{arabic}
\tableofcontents

% ============================================================ 1
\chapter{Principles}

\section{What a WAC test must prove}

\wac{} produces JSON files. That they are valid proves almost nothing: a field
read at the wrong offset, a date interpreted in the wrong time zone, a name
decoded with the wrong code page all give perfectly valid JSON — and wrong
JSON. The costliest defects found in \wac{} were all of that kind: they showed
only when a value was confronted with another source.

The tests therefore fall into three families, which do not answer the same
question:

\begin{description}[leftmargin=0pt,itemsep=0.5em]
  \item[The virtual machine cycle] (chapter~\ref{ch:vm}): does \wac{} run to
    the end on a real Windows, and are its outputs consistent with each other?
    The \emph{cross-checks} confront a value computed by \wac{} with an
    independent value from the same collection.
  \item[The format harnesses] (chapter~\ref{ch:format}): does a decoder give
    exactly the result of the reference — Microsoft's own tool, a published
    standard, another implementation?
  \item[The robustness harnesses] (chapter~\ref{ch:robustesse}): does a parser
    read outside its buffer, or loop for ever, on a truncated or forged file?
\end{description}

\section{Three rules}

\begin{enumerate}
  \item \textbf{A fix comes with the test that would have caught the defect.}
    And that test is checked both ways: it passes on the fixed output, and it
    fails when the defect is replayed.
  \item \textbf{After any change, compare with the previous run}, not just the
    green marks: an artefact that went from 1,522 entries down to 573 lights no
    warning at all (section~\ref{sec:comparer}).
  \item \textbf{Real data does not leave the test workstation.} A real
    collection holds personal data: it is processed locally, its temporary
    copies are deleted after use, and none of it enters the repository
    (\fichier{vmtest/results/} is ignored by git).
\end{enumerate}

% ============================================================ 2
\chapter{Setting up}

\section{The Linux workstation}

Everything is built on Linux, cross-compiled with MinGW-w64; the harnesses run
natively, under \texttt{wine}, or in the virtual machine.

\begin{lstlisting}
sudo apt install g++-mingw-w64-x86-64 wine64 python3 python3-pip
pip install olefile          # independent OLE reader, for the cross-checks
pip install dissect.esedb    # independent ESE reader: WebCache (check skipped otherwise)
# or, leaving the system's Python alone: vmtest/tools/ese-venv
python3 -m venv vmtest/tools/ese-venv && vmtest/tools/ese-venv/bin/pip install dissect.esedb==3.18
\end{lstlisting}

\section{Building WAC and the harnesses}

\begin{lstlisting}
./build-windows.sh --test
\end{lstlisting}

\begin{longtable}{@{}L{7.4cm}L{7.6cm}@{}}
\toprule
\textbf{Output, in \fichier{build-windows/}} & \textbf{Role} \\
\midrule
\endhead
\fichier{WAC.exe} & the collection tool \\
\fichier{raw\_hive\_test.exe} & raw NTFS reader, used by the VM cycle \\
\fichier{tests/lnk\_test.exe} & robustness of the shortcut parser \\
\fichier{tests/parsers\_test.exe} & robustness of the jump list and Prefetch parsers \\
\fichier{tests/evtx\_test.exe}, \fichier{consigne\_test.exe} & event log decoder, exhibit store (chapter~\ref{ch:format}) \\
\fichier{tests/offline\_registry\_test.exe} & hive reader confronted with Microsoft's \\
\fichier{tests/system\_conversions\_test.exe} & WAC's conversions confronted with Windows' \\
\fichier{tests/ntfs\_fixup\_test.exe} & check of NTFS records (torn records) \\
\fichier{tests/ntfs\_runs\_test.exe} & NTFS data runs (a forged list refused) \\
\fichier{tests/deleted\_files\_test.exe} & recovery of deleted files on a volume in memory \\
\fichier{tests/carving\_test.exe} & exact size of a carved PE, ZIP, PDF \\
\fichier{tests/mft\_test.exe}, \fichier{usn\_journal\_test.exe} & \fichier{\$MFT} and change journal against ntfs-3g and the compiler's structures \\
\bottomrule
\end{longtable}

\begin{attention}
The harnesses are deliberately in \fichier{build-windows/tests/}, away from
\fichier{WAC.exe}: they are test tools, with no business on a collection
medium. Never copy the \fichier{tests/} folder onto one.
\end{attention}

The build must end without a single warning: \texttt{-Wall -Wextra} are on, and
warnings have already pointed at real defects.

The other harnesses (chapter~\ref{ch:format}) are each built with a command given
in their header and repeated in this guide; several compile natively on Linux.

\section{The test virtual machine}

The VM is a Windows~11 driven through the \emph{qemu-guest-agent}: no
interaction, no network, commands run as
\texttt{NT AUTHORITY\textbackslash SYSTEM}. Its creation and prerequisites are
described in \fichier{vmtest/README.md}; in short:

\begin{lstlisting}
cd vmtest
ISO_WIN=~/Downloads/Win11_25H2_English_x64.iso ./create-vm.sh   # once
./install-browsers.sh                          # Chrome and Firefox, for web_client (once)
python3 lib/qga.py ping                        # does the agent answer?
python3 lib/qga.py run --shell "dir C:\\"      # run a command
python3 lib/qga.py write <host> <guest>        # send a file
python3 lib/qga.py read  <guest> <host>        # fetch a file
\end{lstlisting}

\begin{note}
A command started with \texttt{qga.py run} is cut after 30 minutes. For a longer
test (for instance the Prefetch files of a real machine,
section~\ref{sec:parsers}), start it as a detached task in the VM and follow its
log.
\end{note}

% ============================================================ 3
\chapter{The test cycle in the virtual machine}\label{ch:vm}

\section{Running a cycle}

\begin{lstlisting}
cd vmtest
./run-wac-test.sh --build     # rebuild, then full cycle
./run-wac-test.sh             # full cycle with the binaries already built
./run-wac-test.sh --raw-only  # only the validation of the raw reading
./run-wac-test.sh --no-split  # without step 7 (collection then conversion)
./run-wac-test.sh --binary-all  # step 7 with --binary-all
./run-wac-test.sh --only=mft,usn_journal  # only these artefacts
\end{lstlisting}

\opt{only} takes keys of the \texttt{artefacts:} section of \fichier{wac.yml}:
the harness places next to \fichier{WAC.exe} a \fichier{wac.yml} that switches
every other one off, for a cycle a few minutes long when a single artefact
changed. The checks of an artefact switched off are skipped (its JSON says
\texttt{NotRequested}); the artefacts read in the hives also need
\texttt{registry}, and step 8 (\fichier{wac.yml}) is skipped. A full cycle
remains the rule before a commit that touches the core.

A second fixed disk of 2~GB (\fichier{create-vm.sh}) is formatted NTFS on the
first cycle, label \texttt{WACDATA}: it holds the \path|$MFT| and journal of a
volume other than the system's, with their witnesses.

The cycle chains: upload of \fichier{WAC.exe} and \fichier{raw\_hive\_test.exe}
into the VM, validation of the raw extraction of a hive (\texttt{reg load} must
refuse it until it is repaired, then accept it), a collection by \wac{} then the
conversion of that collection — \wac{} never does both in the same run —, then
fetch, into
the folder \fichier{vmtest/results/<timestamp>/}, of:

\begin{itemize}
  \item \fichier{run.log} — the console output of the collection, which ends
    with \texttt{END, Time elapsed} when \wac{} went to the end;
  \item \fichier{raw.log} — the trace of the NTFS reader;
  \item the JSON files of the artefacts;
  \item \fichier{exhibits/} — the exhibit manifest, its seal and the list of
    the files really present in the exhibit store.
\end{itemize}

Step 7 then separates collection and conversion (chapter ``Collect here,
convert elsewhere'' of the user documentation): \opt{collect} \opt{events}
\opt{binary}, which must convert nothing; a conversion of an exhibit store with
one retouched exhibit, which must be refused before any read, with its reason;
a conversion, whose JSON files pass the checks below; a second conversion,
whose JSON files must be identical to the first's, byte for byte. Results in
\fichier{split/} and \fichier{split-again/}.

Around this collection, the cycle checks that the reading is
\textbf{harmless}: \fichier{vmtest/guest/harmlessness.ps1} records, before and after,
the three dates and the USN number of two witness files that \opt{binary}
collects (a macro document, a script whose dates are forged to 2001). They must
be identical. A read of the script through Windows (\texttt{Get-Content}) must
then change its access date: otherwise the check would prove nothing about the
dates. The records are in \fichier{harmlessness.txt}; the approach is in the
\emph{Harmlessness file}. Before this collection, the cycle waits 90~s after
its \texttt{taskkill}: Windows writes a registry change to the disk a minute or
so after making it, and the BAM check looks for it in the hive read from the
disk.

\subsection*{Windows XP}

XP has neither guest agent nor PowerShell: \fichier{run-wac-test-xp.sh} drives
it by the keyboard and its serial port.

\begin{lstlisting}
./run-wac-test-xp.sh winxp-test      # XP SP1; winxpsp3-test for SP3
\end{lstlisting}

A CD image carries \fichier{WAC32.exe}, the \fichier{wac.yml} of \fichier{bin/}
and \fichier{run-xp.cmd}, which the script types into the Run box: it
collects, converts, says each step and its exit code on \texttt{COM1},
deletes \fichier{WAC32.exe} and shuts the machine down. The JSON, the
manifest and the logs are then read on the disk of the VM, off, read-only
(\fichier{lib/read-offline.py}, dissect), and checked by
\fichier{check-json.py}. Without PowerShell no witness is placed: the checks
bear on the outputs' consistency and on what the version must have or not
(IPv4-only network tables, Prefetch 17 without \path|$MFT| reference,
\fichier{Thumbs.db} without checksum). \fichier{run-xp.cmd} spans several
lines on purpose: it checks that \wac{} gives the console its code page
back, without which the \fichier{cmd.exe} of before Windows~8 stops reading a
batch file after it.

\section{Checking the outputs}

\begin{lstlisting}
python3 lib/check-json.py results/<timestamp>
\end{lstlisting}

Each line starts with a green tick (check passed), a red cross (defect) or a
``skip'' sign (check not applicable: the artefact is absent or empty). The exit
code is 0 if no check failed. Beyond JSON validity and the shape of paths, the
script confronts independent sources:

\begin{longtable}{@{}L{6.6cm}L{8.4cm}@{}}
\toprule
\textbf{Cross-check} & \textbf{What it has already caught} \\
\midrule
\endhead
Prefetch hash / suffix of its file name & 65 wrong hashes out of 270 (hexadecimal without padding) \\
Prefetch executable name / start of the file name & a name read 8 bytes too early (``\textbackslash x11'' everywhere) \\
Prefetch \texttt{\textbackslash VOLUME\{guid\}} paths resolved & 20,066 paths out of 20,066 unresolved \\
\texttt{X} / \texttt{XUtc} pairs, lists included & double time-zone shift; Prefetch \texttt{RunsUtc} 2~h off \\
offset of every local date / tz database & the collection day's offset applied to every date; daylight saving dates misread \\
no session earlier than the boot & date shifts, a boot time 3.5~s late \\
ShimCache date / NTFS date in the manifest & a UTC date treated as local (-2~h) \\
DestList entries / jump list shortcuts & the first 15 entries of every list lost \\
machine name of the events / \fichier{OperatingSystem.json} & check of the BinXML decoding \\
record identifiers unique per log & a stale chunk read again \\
mounts / \texttt{reg query} and \texttt{Get-Partition} & device paths published as ``\textbackslash{}'' \\
precision of FAT and text dates & seven digits of fraction on dates to the second \\
BAM, UserAssist, USBSTOR, Amcache / harness, Prefetch, Windows, PE headers & UTC dates read as local (2~h) \\
SID names / \texttt{Get-LocalUser}, \texttt{Get-LocalGroup} & names asked of the system, translated (``Système'') \\
GUIDs registered by the harness (\fichier{register-test-guids.ps1}) & ``Unmapped GUID'' published as a name \\
seal / real fingerprint of the manifest & a retouched manifest \\
files in the exhibit store / manifest content & 209 exhibits added after sealing \\
SHA-256 of files read raw / \texttt{Get-FileHash} & safeguard of the reading by batches of clusters \\
Prefetch dates / dates given by Windows & 0 out of 335 matching: dates of the working copy \\
executables of System32 and SysWOW64 / manifest of a \opt{collect} \opt{binary} & 6,446 missing out of 8,540 (\texttt{\$INDEX\_ROOT} in an extension record, hard links) \\
catalogs cited in a verdict / exhibits of the store & 170 out of 170 missing (label ``catalog'' looked for as ``catalogue'') \\
\wac{}'s Microsoft verdicts / \texttt{Get-AuthenticodeSignature} & safeguard: a wrong verdict would leave a binary out of the store \\
copy of a Store package (untouched, modified, added) / verdicts & verification through the signed block map \\
two conversions of one collection & dates read on the working copy, rebuilt at each conversion \\
forged dates of a script / \path|$STANDARD_INFORMATION| and \path|$FILE_NAME| & safeguard: the forgery must stay visible \\
retrieval keys / PE headers read by PowerShell & safeguard of the retrievable content \\
\opt{binary-all} / copy and verdict of every executable & package check too strict under \opt{binary-all} \\
\fichier{verify-exhibits.ps1} / separate collection & fingerprint-only exhibits reported altered \\
\bottomrule
\end{longtable}

The full list, with the history of each check, is in \fichier{vmtest/README.md}.

\section{Comparing two runs}\label{sec:comparer}

A green mark does not tell whether a piece of data disappeared. After any
change, compare the new collection with the previous one: number of entries of
each artefact, and values lost or gained. Two precautions:

\begin{itemize}
  \item a collection changes by itself from one run to the next: timestamps,
    new events, updated Prefetch files, the fingerprint of \fichier{\$MFT}.
    Compare \emph{paired} items — for instance a Prefetch by its path hash and
    its last run — leaving aside the dates of the extracted copy;
  \item a value that looks ``lost'' may just be an empty field no longer emitted, or a
    renamed key: normalise before concluding.
\end{itemize}

% ============================================================ 4
\chapter{The format harnesses}\label{ch:format}

Each one confronts a \wac{} decoder with a reference. All are excluded from the
\wac{} build by their \fichier{\_test.cpp} suffix.

\section{Fingerprints: \texttt{sha\_test}}

Confronts SHA-1 and SHA-256 with the published vectors of FIPS~180-4, and with
lengths 55 to 128 bytes, which exercise the padding — the one place where an
otherwise correct implementation goes wrong.

\begin{lstlisting}
cd sources && g++ -std=c++17 -I. sha.cpp sha_test.cpp -o /tmp/sha_test && /tmp/sha_test
\end{lstlisting}

\section{JSON reading: \texttt{json\_test}}

\texttt{--convert} reads back what \texttt{--collect} wrote: the exhibit
manifest and the snapshot of the live state. The harness checks that
everything WAC's writer produces reads back identical, that the files of a real
collection read back and rewrite to the byte, that malformed texts are refused,
and that 200,000 random mutations never make the reader read out of bounds
(AddressSanitizer, UBSan).

\begin{lstlisting}
cd sources && g++ -std=c++17 -g -fsanitize=address,undefined -I. json_test.cpp -o /tmp/json_test
/tmp/json_test ../vmtest/results/<timestamp>/*.json
\end{lstlisting}

\section{Decompression: \texttt{lznt1\_test} and \texttt{xpress\_test}}

Check the NTFS (LZNT1) and WOF (XPRESS Huffman) decompressors against
\emph{Microsoft's own compressor}: a known file is compressed in the VM by
\texttt{RtlCompressBuffer}, unit by unit, and the harness decompresses each unit
and compares it byte for byte with the original. A wrong decompressor usually
decodes without any error: only the exhaustive comparison reveals it.

\begin{lstlisting}
cd WAC
g++ -std=c++17 -I. lznt1.cpp lznt1_test.cpp -o /tmp/lznt1_test
g++ -std=c++17 -I. xpress.cpp xpress_test.cpp -o /tmp/xpress_test
/tmp/lznt1_test  <original> <compressed.bin> <index.idx>
/tmp/xpress_test <original> <compressed.bin> <index.idx>
\end{lstlisting}

Each unit has one line in \fichier{index.idx}: \texttt{<n> <original size>
<compressed size>}; a null compressed size marks a unit Windows gave up
compressing.

\section{Artefacts by version of Windows: \texttt{artefact\_catalogue\_test}}

Checks the table that says, for each version of Windows, which artefacts
exist, where and in which format (\fichier{artefact\_catalogue.cpp}). That table
decides what a collection does \emph{not} take: a wrong line would lose an
exhibit without a word. The test checks that the table and \fichier{wac.yml}
hold exactly the same keys, that each artefact's forms are in order, that known
cases answer right (no SRUM before Windows~8, Prefetch~23 on Vista and 30 on
Windows~10), that on Windows~7 the guard holds back only the five artefacts
that version lacks, and that a Windows newer than the table is treated as the
last one known, and flagged. An unknown key or version never holds a collection
back. Only a server edition is said not validated by the test cycle: every
client, from XP SP1 to Windows~11, is. It also checks that a source is said ``always there'' only where it is:
not on a server, not for a key that exists only once a device was plugged, not
for a metadata file of the volume or a state of the running system. Lastly,
that a piece is decoded only if the table holds its signature with a decoder: a
Prefetch of version~23 is named and not decoded, a Prefetch of version~32,
which nobody has seen yet, has neither form nor decoder. Seen
failing: SRUM declared present since XP; \fichier{USBSTOR} declared always
there.

\begin{lstlisting}
cd sources && g++ -std=c++17 -g -fsanitize=address,undefined -I. artefact_catalogue.cpp artefact_catalogue_test.cpp -o /tmp/act
/tmp/act wac.yml
/tmp/act wac.yml --dump    # the table, for doc/artefact-table.py
\end{lstlisting}

In the cycle, the version \wac{} reads in the copied \fichier{SOFTWARE} hive is
confronted with the one Windows gives of itself (build number, workstation or
server).

The last step of the cycle moves the \fichier{hosts} file aside for the time
of a collection of that artefact alone, then puts it back. The conversion must
give \fichier{hosts.json} the status \texttt{ExpectedButMissing}, with the
build Windows reports and the place where the file should have been, and record
the finding. The independent reference is the manifest of the exhibit store: it
holds the other files of the folder --- the folder was read --- and no
\fichier{hosts}. No other artefact may carry the status; in the rest of the
cycle, where the test machine lacks nothing, it is a fault. Seen failing on an
ordinary collection.

A format without decoder is presented to the decoders themselves.
\texttt{parsers\_test} rewrites the version of a real Prefetch to the next
one: it must be declared unsupported, the name and the hash
read, nothing else decoded. \texttt{file\_artefacts\_test} does the same for
a \fichier{\$I} index of version~3, a thumbnail cache of Vista and a jump
list DestList of version 5 or 7 (version~6 of Windows~11, measured with
\texttt{olefile}, is read). In the
cycle, a Prefetch file written by hand (version 32, never seen) is placed in
\path|C:\Windows\Prefetch|: \fichier{prefetchs.json} must declare it without
naming a format, and declare no real Prefetch --- every version from XP's on
is decoded. Every real Prefetch of the machine is then read by libscca, written
independently of \wac{}: version, hash, run count, run times to the 100~ns,
volume serials and number of files must be its.

\section{Browser sessions: \texttt{browser\_sessions\_test}}

Checks the decoding of Chromium's session files (SNSS: \path|Session_*|,
\path|Tabs_*|) and of Firefox's mozLz4 files (\fichier{.jsonlz4}). The LZ4
blocks are made by Python's \texttt{lz4} module, an independent
implementation: each must decompress byte for byte. The SNSS files are written
by the test field by field, as Chromium lays them out (Pickle, fields aligned
to 32 bits): a navigation replaced by a later one, an older navigation without
the fields after its title, a closed tab and its instant, a truncated command,
a lying length. Then every input, damaged at random (fixed seed), must never
make the decoders read out of bounds (AddressSanitizer, UBSan). Seen failing:
without the fields' alignment, and with an LZ4 copy that does not overlap.

\begin{lstlisting}
cd sources && g++ -std=c++17 -g -fsanitize=address,undefined -I. browser_sessions.cpp browser_sessions_test.cpp -o /tmp/bst
python3 -c "import lz4.block,sys; open(sys.argv[2],'wb').write(lz4.block.compress(open(sys.argv[1],'rb').read(), store_size=False))" a.bin a.lz4
/tmp/bst a.bin a.lz4 [<original> <block> ...]
\end{lstlisting}

In the cycle, Chromium's sessions are confronted with \fichier{History}: each
page the history knows carries the instant of one of its visits, to the
microsecond; Firefox's, decompressed by Python's \texttt{lz4} module, must give
the same pages, states and closing instants.

\section{LZX decompression: \texttt{lzx\_test}}

WOF's fourth algorithm, LZX on 32~KiB chunks, is only produced on request. The
test data come from Windows itself: in the VM, \texttt{compact /c /exe:lzx}
compresses a file, and \texttt{raw\_hive\_test --wof} writes its compressed
stream as it lies on the disk. The test decompresses every chunk and compares it
with the original, then submits every chunk, truncated and damaged at random
(fixed seed), to the decoder built with AddressSanitizer and UBSan: no read or
write out of bounds, and never a ``success'' of the wrong size.

\begin{lstlisting}
cd WAC
g++ -std=c++17 -g -fsanitize=address,undefined -I. lzx.cpp lzx_test.cpp -o /tmp/lzx_test
/tmp/lzx_test <original> <wof-stream.bin>
\end{lstlisting}

Reference: \fichier{ntoskrnl.exe}, \fichier{shell32.dll} and a text file,
646 chunks out of 646 identical, 235,707 checks, no failure. The VM cycle also
checks, on every run, the SHA-256 of these three files compressed with LZX and
read raw against \texttt{Get-FileHash}.

\section{Event messages: \texttt{wevt\_test} and \texttt{event\_messages\_test}}

\texttt{wevt\_test} checks the chain that turns an event into a sentence, on the
files of a real provider: its DLL and its localised satellite.

\begin{lstlisting}
cd sources && g++ -std=c++17 -I. pe_resource.cpp wevt.cpp wevt_test.cpp -o /tmp/wevt_test
/tmp/wevt_test gpsvc.dll fr-FR/gpsvc.dll.mui "{aea1b4fa-97d1-45f2-a64c-4d69fffd92c9}" 1002 4001
\end{lstlisting}

\texttt{event\_messages\_test} does the same starting from the SOFTWARE hive; it
runs on the examined machine or in the VM, the only places where the hive and
the providers' binaries are readable:

\begin{lstlisting}
event_messages_test <SOFTWARE hive> <guid> [id[:version] ...]
\end{lstlisting}

\section{Event logs: \texttt{evtx\_test}}

Reads an \fichier{.evtx} outside Windows, under \texttt{wine}; built by
\texttt{build-windows.sh --test}, in \fichier{build-windows/tests/}. It serves to try
the decoder on logs the VM does not produce: old, forwarded, closed abruptly,
damaged — and to compare record by record with an independent implementation
(\texttt{python-evtx}).

\begin{lstlisting}
wine evtx_test.exe "Z:/path/Security.evtx"            # summary per log
wine evtx_test.exe "Z:/path/Security.evtx" 3          # + XML of 3 records
wine evtx_test.exe "Z:/path/Security.evtx" --dump     # one record per line
wine evtx_test.exe --collect "Z:/tmp/tree" "Z:/tmp/out"          # complete chain
wine evtx_test.exe --collect-memory "Z:/tmp/tree" "Z:/tmp/out2"  # same chain, in memory
cmp /tmp/out/events.json /tmp/out2/events.json                   # must be silent
\end{lstlisting}


\section{Hive reader: \texttt{offline\_registry\_test}}

\wac{} reads registry hives with its own reader
(\fichier{offline\_registry.cpp}), linked into the executable: it no longer
depends on Microsoft's \fichier{offreg.dll}, which Windows ships and which it
therefore borrowed from the examined machine. This harness confronts the two:
it loads Microsoft's DLL from the path given, walks each hive with both readers
and compares every key (name, date, class, counts, maxima, security descriptor
size) and every value (name, type, bytes), as well as the contract's edge
cases: buffer too small, size probe, default value, missing key or value. No
difference is allowed.

\begin{lstlisting}
offline_registry_test.exe C:\Windows\System32\offreg.dll SYSTEM SOFTWARE SAM ntuser.dat ...
\end{lstlisting}

To be run on Windows (the VM), on replayed hives: those of a collection's
\fichier{working/} folder. The one intended difference: a hive whose logs were
not replayed is refused by \wac{}, where recent versions of Microsoft's DLL
read it as is — stale keys, without an error; the harness reports it as
expected. Reference: 18 hives, ten of them from a real machine, i.e. 1,273,652
keys and 2,456,139 values, identical.

\section{Conversions: \texttt{system\_conversions\_test}}

\wac{} imports only what querying the running system requires
(\fichier{ADVAPI32}, \fichier{KERNEL32}, \fichier{msvcrt}, \fichier{Secur32},
\fichier{WTSAPI32}). The conversions it used to ask of other libraries are
done by its own code; this harness confronts them with the Windows functions
they replaced: GUIDs (\texttt{StringFromGUID2}), SIDs both ways
(\texttt{ConvertSidToStringSidW}, \texttt{ConvertStringSidToSidW}), OLE dates
(\texttt{VariantTimeToSystemTime}) and property names
(\texttt{PSGetNameFromPropertyKey}), FAT dates (\texttt{DosDateTimeToFileTime},
on every date), and daylight saving time: WAC's offset in both directions
(\texttt{SystemTimeToTzSpecificLocalTimeEx}, \texttt{TzSpecificLocalTimeToSystemTimeEx})
on \emph{every} time zone of Windows, with its rules year by year, at the
transitions (hour skipped, hour repeated, change of rule on 1 January) from
1980 to 2035; the reading of \fichier{SYSTEM\textbackslash...\textbackslash
TimeZoneInformation}, whose daylight saving dates are \texttt{TIME\_FIELDS},
not \texttt{SYSTEMTIME}s; rules forged at random; the reading back of the
manifest's ISO~8601 dates (\texttt{iso8601UtcToFiletime}, exact to the 100~ns on
200,000 instants, malformed dates refused). It also checks the
suspect's offset in both directions: a local
date converted to UTC must be so with the time zone of the examined machine,
not that of the machine running \wac{}. An impossible date (month 13, hour 25)
gives a null date, hence not emitted. The table of property names is generated
from the schema of a reference Windows by \fichier{vmtest/guest/property-schema.cpp}.

\begin{lstlisting}
system_conversions_test.exe
\end{lstlisting}

To be run on Windows (the VM). Reference: 9,861,570 comparisons, no difference.
Two intended differences: \wac{} refuses a SID text Windows would truncate (a
sub-authority of $2^{32}$ or more, a sixteenth sub-authority), and a NaN date,
for which Windows returns a meaningless time.

\section{Exhibit store: \texttt{consigne\_test}}

Replays the production chain — fingerprints, manifest, verified copy, hive
replay — on hives given as arguments, and checks what no build reveals: the
exhibit store stays identical byte for byte while the working copy is changed,
and the seal carries the manifest's real fingerprint. Last, it checks that an
exhibit path read from a manifest stays inside the exhibit store: a
\texttt{..} component is refused, a WinSxS name holding \texttt{..}
(\texttt{i..sermode}) is accepted. An operation written to the investigation log
before the suspect's time zone is read still carries its local time, never
UTC labelled \texttt{+00:00}; and an exhibit store reloaded as a conversion
does gives its exhibits' size, one byte too little free space being refused
for the space, not as invalid data.

\begin{lstlisting}
cd build-windows/tests
wine consigne_test.exe "Z:/tmp/out" "Z:/path/SYSTEM" "Z:/path/ntuser.dat"
\end{lstlisting}

\section{Replay of the hive logs: \texttt{hive\_recover\_test}}

A replay that drops an entry, or applies a stale one, gives a hive that still
opens and still reads: only the keys of the lost pages are missing, without any
error. The test forges a hive and its logs where each page holds one letter and
each entry writes another: the letter left in a page says which entry was
applied there last. It checks that the entry numbered as the hive's sequence is
applied — it used to be skipped, one entry lost at every replay —, that the
earlier entries and the leftovers of a previous generation are not, that an
entry growing the hive extends the file and the header, that the header ends
clean with its checksum, that the undo log gives the original hive back, and
that an altered entry ends the chain.

\begin{lstlisting}
cd build-windows/tests
wine hive_recover_test.exe "Z:/tmp/work"
\end{lstlisting}

Reference: 11 checks, no failure; four fail if the comparison of the sequence
numbers is put back to its former form. On real hives, the VM cycle confronts
\wac{}'s replay with Windows' own recovery of the same files (keys, values and
a digest of them all).

\section{Hashed key of the SAM: \texttt{sam\_secrets\_test}}

A wrong constant in the SAM's RC4 scheme (Windows~8.1 and earlier, and the
accounts nothing has rewritten since) gives a wrong hashed boot key, then
wrong NT hashes for every account, without any error. The test builds an
\texttt{F} value from a synthetic boot key and hashed boot key, encrypted and
given its checksum by an independent computation; it checks that the hashed
boot key comes back, and that a wrong boot key, an altered checksum, a
truncated value or an unknown revision are refused: no key rather than a
wrong one.

\begin{lstlisting}
cd build-windows/tests
wine sam_secrets_test.exe
\end{lstlisting}

Reference: 5 checks, no failure; the first fails with the former constant
(thirteen ``Q'' instead of twelve). On a real SAM, the VM cycle confronts a
witness account's hash with MD4 of its password, on Windows~11 (AES) and on
Windows~8.1 (RC4).

\section{Shadow copies read raw: \texttt{vss\_test}}

A wrong rule in the application of the shadow copy stores gives a copy whose
files are valid and wrong: a block of the current volume where the copy kept
another. Without argument, the test builds a volume of 64 blocks and two copies
whose answers are set by hand, block by block: a block the older store keeps, a
forwarder to the newer store, a forwarder taken over by a later descriptor, an
overlay of one sector that serves its own copy only, an ignored descriptor, an
unused block read as zeros for the newer copy alone, a catalog given out of
order, a block Windows copies after the lists were read (seen after a refresh,
the older store kept), a catalog chain that loops (refused).

\begin{lstlisting}
cd build-windows/tests
wine vss_test.exe
wine vss_test.exe Z:\path\volume.raw   # image of a volume: list and SHA-256 of each copy
\end{lstlisting}

Reference: 30 checks, no failure; seen failing on four defects put back
(overlay applied to every copy, unused blocks not zeroed, forwarder chain not
merged, older stores lost by the refresh). On a real 2~GiB volume holding two
copies, the list and the SHA-256 of each copy's rebuilt volume are
libvshadow's; the witness files Windows reads through the copy's device are the
same. The VM cycle confronts the manifest's fingerprints with those Windows
computes through the copy the harness made.

\section{Verification scripts: \texttt{verify\_exhibits\_test.sh}}

Runs \fichier{verify-exhibits.sh} and \fichier{verifier-consigne.sh}, the
Linux scripts delivered with \fichier{WAC.exe}, on forged collections whose
verdict is known: an intact collection (a path with a space, a fingerprint in
lower case, a fingerprint-only binary) is accepted; an altered, missing or
unfingerprinted exhibit, a modified manifest, an unreadable one and a sealed
one without its list of exhibits (or with no list) are refused. The seal
written down at the collection (\texttt{-{}-seal}) accepts the collection it
sealed, in either case, and refuses a manifest rewritten and resealed.
A manifest that names a file of the investigation machine — by \texttt{..}, by
an absolute path or by a symbolic link — is refused without the file being
read. \fichier{verifier-consigne.sh} being the same script in French, every
case runs through both, with the same exit code and counters expected. Linux only, no data
needed.

\begin{lstlisting}
bash sources/verify_exhibits_test.sh
\end{lstlisting}

Reference: 23 checks, no failure. The \fichier{.ps1} scripts are run by the
VM cycle on the separate collection.

\section{Raw reading: \texttt{raw\_hive\_test}}

Extracts a file by reading the NTFS volume directly. To be run on Windows, as an
administrator; the VM cycle uses it to validate the extraction of a hive.

\begin{lstlisting}
raw_hive_test.exe [volume] [path] [output] [--fix] [--list] [--attrs]
\end{lstlisting}

\texttt{--fix} makes the extracted hive loadable, \texttt{--list} lists a
directory, \texttt{--attrs} shows the \fichier{\$MFT} attributes of the file
without interpreting them.

\section{NTFS records: \texttt{ntfs\_fixup\_test}}

NTFS writes a \fichier{\$MFT} record or an index block sector by sector,
and ends every 512-byte stride with a sequence number. \wac{} checks that
number before restoring the real bytes: a record torn by an interrupted write
is refused. The harness builds records whose expected result is known (sound,
torn in each stride, malformed arrays), then confronts \wac{}'s verdict with an
independent check on 20,000 random records placed against a forbidden page.
No data is needed.

\begin{lstlisting}
cd build-windows/tests && wine ntfs_fixup_test.exe
\end{lstlisting}

Reference: 44,610 checks, no failure. On a real volume, the VM cycle confirms
it: \fichier{raw.log} must hold no ``torn record''.

\section{NTFS data runs: \texttt{ntfs\_runs\_test}}

The data runs say where a file's clusters lie; read from a deleted record,
they are what an intruder or a damaged disk left. The harness decodes known
lists (the Linux-NTFS documentation's examples, a run backwards, a sparse run,
the attribute's end without a zero), refuses forged lists whole (a size field
beyond 8~bytes, a length of zero, a truncated run, a cluster before the
volume's start), then confronts \wac{} with an independent decoder on 128-bit
integers for 50,000 random lists placed against a forbidden page.

\begin{lstlisting}
cd build-windows/tests && wine ntfs_runs_test.exe
\end{lstlisting}

Reference: 100,026 checks, no failure; the former decoder, which took a field
of 15~bytes, fails 48,224 of them.

\section{Deleted files: \texttt{deleted\_files\_test}}

A recovery gone wrong gives a file that looks right: the bytes of whatever
lies in its clusters now. The harness builds deleted records and a volume in
memory whose outcome is known: a resident file, a contiguous one, a fragmented
one with a sparse run; a cluster taken before the reading, another taken
\emph{during} it; all taken; clusters read back zero (TRIM); a valid data
length shorter than the size. Inconsistent lists (a run outside the volume,
two runs on the same clusters, sizes that do not match, compression,
encryption, runs in an extension record) are refused before any read. Every
read asked for is checked: inside the volume, on free clusters. Then 20,000
damaged records, against a forbidden page. Level~2: an INDX block built with
live entries and, in its slack, one valid entry among decoys (another parent
directory, an 8.3 alias, a key of the wrong length, a date of 1917, a control
character) — only the valid one comes out; a released block also gives its
former live entries; 4,000 damaged blocks against a forbidden page.

\begin{lstlisting}
cd build-windows/tests && wine deleted_files_test.exe
\end{lstlisting}

Reference: 91,842 checks, no failure — the clusters claimed since by a later deleted file among them (known cases, 2,000 draws against a cluster-by-cluster reckoning, a hole said so); without the check of the parent
directory, level~2 fails three of them; without the check of \fichier{\$Bitmap}
after the reading, the cluster taken during the reading is published as the
file's content and the check fails. In the VM, files of known content
(resident, text, unsigned executable, \fichier{notepad.exe}) are written to
the disk then deleted: \texttt{check-json.py} requires a content declared
\texttt{Recovered} to have the SHA-256 Windows computed before the deletion,
and the authentic executable to be listed without being copied; a fifth witness is written over after its deletion, by a file deleted in turn: never ``Recovered'' with another content. For level~2,
150 files are moved from a folder to another: at least one former name must
come out of the slack of the first one's index, with the record and current
path \fichier{mft.json} gives the file.

\section{Carving: \texttt{carving\_test}}

A carved file has only its bytes: cut one byte short or long, it looks like
the file and is not. The harness builds a PE (two sections, a signature), a
ZIP (one file, a comment) and a PDF (an xref table) of known size, and their
forged variants (unknown machine, 200 sections, a section off its alignment,
a signature inside the sections, a central directory elsewhere than the end
record says, a \texttt{startxref} pointing at nothing); then 30,000 damaged
heads against a forbidden page; then real files given as arguments, which
must come out at their exact size.

\begin{lstlisting}
cd build-windows/tests && wine carving_test.exe 'Z:\...\WAC.exe' 'Z:\...\documentation_FR.pdf' 'Z:\...\kape.zip'
\end{lstlisting}

Reference: 96,760 checks, no failure; WAC.exe and WAC32.exe (with their COFF
symbol table), a PDF manual, an archive of 140~MB to the byte. In the VM
(\texttt{WAC\_DELETED\_LEVEL=3}, XP and Windows~11), an executable and a PDF
written, truncated to nothing then deleted must be carved with the SHA-256 of
their writing; on the FAT32 and exFAT keys, a deleted text must be read back
with its own.

\section{\fichier{\$MFT} and change journal: \texttt{mft\_test}, \texttt{usn\_journal\_test}}

An offset error in reading the \fichier{\$MFT} or the change journal gives
dates, names and references that are valid and false: no syntax check would
see it. Both harnesses therefore compare \wac{} to an independent reference.

\fichier{mft\_fixtures.sh} formats an NTFS volume in a file with
\texttt{mkntfs}, drops files in it (resident, non-resident, a non-ASCII name)
with \texttt{ntfscp}, without mounting or privilege, dumps its
\fichier{\$MFT} with \texttt{ntfscat}, then has each record described by
\texttt{ntfsinfo} — ntfs-3g, another implementation of NTFS.
\texttt{mft\_test} compares, record by record, the name, the parent, the
size, the eight dates and the named streams read by \wac{} to that
description. \texttt{usn\_journal\_test} builds records of versions~2 and~3
with structures whose layout comes from the compiler (MinGW's
\texttt{USN\_RECORD}; version~3, which MinGW does not declare, from
Microsoft's documentation, its offsets checked at compile time), lays them out
as a journal the way NTFS leaves it (pages, zeros, a range record, a leading
sparse region), checks the numbers, the reasons and the paths, then decodes
every record damaged at random against a no-access page.

\begin{lstlisting}
bash sources/mft_fixtures.sh /tmp/mft          # ntfs-3g: mkntfs, ntfscp, ntfsinfo
cd build-windows/tests
wine mft_test.exe "Z:/tmp/mft"
wine usn_journal_test.exe "Z:/tmp/mft"
\end{lstlisting}

Reference: 281 checks for \texttt{mft\_test}, 27 for
\texttt{usn\_journal\_test}, no failure.

\section{Raw reading on a real volume: \texttt{raw\_volume\_test.sh}}

The directory index names a file by its record number \emph{and} a sequence
number, which NTFS increments every time it frees the record. On the running
system a file can be deleted between the listing of its directory and its
read, and its record given to another file: \wac{} used to read it anyway,
and sealed that other file's content under the first one's name (an
\fichier{SRU*.log} holding a hive's \texttt{hbin} blocks, found in the test
VM). The script serves the volume of \fichier{mft\_fixtures.sh} to Wine as
drive \texttt{T:}, in a Wine prefix of its own, and has
\texttt{raw\_hive\_test} read it as \wac{} reads a disk: a file whose bytes
must be those \texttt{ntfscat} reads, then the same file once its record has
another sequence number than its index entry — refused, no copy left.

\begin{lstlisting}
bash sources/mft_fixtures.sh /tmp/mft
bash sources/raw_volume_test.sh /tmp/mft build-windows/raw_hive_test.exe
\end{lstlisting}

Reference: 3 checks, no failure; the reader before the fix fails 2 of them.

\section{Outlook's mail: \texttt{pst\_test}}

A PST read wrong — a block decrypted with the wrong table, a heap item cut one
byte off, a compressed block inflated short — still gives valid JSON.
\fichier{pst\_fixtures.sh} downloads, pinned by their SHA-256, java-libpst's
test files (Apache 2.0: two Unicode PSTs and an OST of Outlook 2013, 4~KiB
pages and compressed blocks) and a PST of the EDRM Enron v2 corpus (CC-BY 3.0
US), and has each message read by pypff, libpff's reader written independently
of \wac{}. \texttt{pst\_test} compares, message by message: subject, sender,
delivery and submission times, the name, size and SHA-256 of every
attachment, the SHA-256 of the body.

\begin{lstlisting}
PYTHON=vmtest/tools/pff-venv/bin/python bash sources/pst_fixtures.sh /tmp/pstfx
cd build-windows/tests && wine pst_test.exe Z:/tmp/pstfx
\end{lstlisting}

Reference: 1,800 checks, no failure; the 256 messages pypff reads identical,
and the 4 items of the Enron PST pypff cannot read read by \wac{}. With the
decryption removed, 253 failures.

\section{Internet Explorer's indexes: \texttt{ie\_index\_dat\_test}}

Checks the decoder of Internet Explorer 5 to 9's \fichier{index.dat} files
(Windows 7 and earlier) on four real files of the plaso project's public test
data — history, cache, daily history, and a cache of over a thousand records
with redirections, dropped entries (\texttt{LEAK}) and recovered ones. The
reference is \texttt{pymsiecf}, libmsiecf's reader, written independently of
\wac{}: every record it gives must come back with the same fields (URL, file,
cache directory, times, expiry, last check, hits, data size), and none more.
This test caught a real defect: a recovered record, partly overwritten by a
newer one, gave the other's data as its own. Then every file is decoded
truncated and damaged at random, against a no-access page on both sides;
finally the conversion (container, account and page, cached file, headers,
times) is checked.

\begin{lstlisting}
pip install libmsiecf-python              # pymsiecf, for the reference
PYTHON=python3 sources/ie_index_fixtures.sh /tmp/iefx
cd build-windows/tests && wine ie_index_dat_test.exe Z:/tmp/iefx
\end{lstlisting}

The files are not versioned: the script downloads them and refuses them if
their SHA-256 differs. Reference: 22 checks, no failure (1,110 records
identical to \texttt{pymsiecf}'s).

\section{Windows Timeline: \texttt{activities\_cache\_test}}

Checks the reading of \fichier{ActivitiesCache.db} on a real Windows~10
database of the plaso project's public test data (112 activities:
applications, pages, foreground time). The reference is Python's
\texttt{sqlite3}, each column converted by the rules of
\fichier{activities\_cache.h} written again in the script, independently of
\wac{}: Unix times in UTC, GUIDs in .NET's byte order, JSON payloads parsed,
null values left out. Every row must come back with the same columns and
values, and none more; the local forms of the dates are not compared (they
depend on the examined machine's time zone). Then a truncated copy is read
without a crash. Reintroducing a one-hour shift in the conversion of the
dates makes 113 checks fail.

\begin{lstlisting}
sources/activities_cache_fixtures.sh /tmp/acfx
cd build-windows/tests && wine activities_cache_test.exe Z:/tmp/acfx
\end{lstlisting}

The database is not versioned: the script downloads it and refuses it if its
SHA-256 differs. Reference: 118 checks, no failure.

\section{Small file artefacts: \texttt{file\_artefacts\_test}}

Checks the decoders of the PowerShell history, of \fichier{setupapi.dev.log},
of the recycle bin's \path|$I|, of \fichier{hosts}, of \fichier{Report.wer}
and of Defender's MPLog on inputs written from each format's description, with
the expected result: a command on two lines joined, a section title split into
operation and subject, 31 February refused as a date, a \path|$I| of version 1
and 2 read, a \path|$I| truncated or of unknown version refused, a path length
that lies bounded by the file, a problem signature ordered by its index
whatever the file's order, a date without \texttt{Z} not taken as UTC. For
\path|$Secure:$SDS|: a stream built block by block, where the mirror block is
skipped even when its entry looks consistent, a repeated \champ{SecurityId}
kept once, a length leaving the block refused, a damaged hash flagged
(\champ{HashMismatch}), the owner and the ACE decoded, every read against a
no-access page on both sides. For \path|$Boot|: geometry and serial read, a
sectors-per-cluster count above 0x80 taken as a power of two, a sector without
0x55AA, of another file system, of 100 bytes or truncated refused; for
\path|$Max|: the four fields, a stream too short refused. For a DPAPI master
key: header read, block sizes beyond the file or wrapping when added refused,
a truncated header refused; \fichier{Preferred} given in the form of the key
files' names. For \path|User Shell Folders|: \texttt{\%USERPROFILE\%}
replaced by the profile whatever its case, a path on another disk kept, a share
or another user variable not followed. For the SQLite reading: a database in
WAL mode whose writer is still open, as a running browser leaves it, read on
its private copy with its log's rows (UTF-8 text, integer, NULL, BLOB), the
database and its log unchanged, the copy deleted, a query of another schema
and a file that is no database refused. For the Chromium history: a
\fichier{History} built with Chromium's tables, a typed visit then a link
citing it, a redirected download whose chain comes back in its order, a
searched term; transition types and qualifiers, Chromium times converted. For
Firefox: \fichier{places.sqlite} and \fichier{formhistory.sqlite} built with
Firefox's tables, a download visit citing the visit it came from, a download's
destination decoded from its \texttt{file:///} URI (UTF-8, a space), its end
in milliseconds, an address bar input and a search bar entry. The VM has no
Firefox: this test is its only check. Saved logins and cookies of Chromium and
Firefox: site, account, dates (Firefox cookies' expiry in seconds), and no
password nor cookie value in the JSON. No data needed.

\begin{lstlisting}
cd build-windows/tests && wine file_artefacts_test.exe
\end{lstlisting}

Reference: 283 checks, no failure (the module list of the test's own
process, its executable first, read as \wac{} reads it; the second-lot registry values included: TrustRecords, firewall rule, certificate Blob; and the thumbnail caches: a real entry header of the VM and the checksum Windows wrote for it, a cache built here, forged lengths; Thunderbird's mail: an mbox written here, encoded words, quoted-printable, base64, an RFC~2231 name, a nested multipart or one without boundary, a date with its offset). In the VM cycle,
\fichier{place-file-artefacts.ps1} places witnesses whose content is known (a
history, a file sent to the bin, a hosts entry, a startup script) and
\fichier{check-json.py} finds them in the JSON.

\section{Signatures: \texttt{authenticode\_test}}

Verifies the catalogs of a folder, then gives a verdict per file; its verdicts
are compared with those of \texttt{Get-AuthenticodeSignature} on the same files,
in the VM. The only gap that matters: a file \wac{} accepts and Windows refuses.

\begin{lstlisting}
authenticode_test --catalogs <folder>
authenticode_test --catalogs <folder> --files <list>
\end{lstlisting}

\section{Configuration: \texttt{yaml\_light\_test}, \texttt{config\_test}}

\texttt{yaml\_light\_test} (native, libyaml included, under AddressSanitizer
and UBSan): exact readings (nested sections, quotes, comments, byte order
mark, empty file), refusals with the right line (duplicate key, list outside
a key, list in a list, mapping in a list, anchor, alias, tag, second
document, tab, quote left open, nesting too deep), lists of values read
exactly (the deleted files' excluded folders), then 3,000 texts damaged at
random; seven mutations all detected.
\texttt{config\_test} (under wine): each setting applied, each kind of mistake
refused with its line, and the reference file of \opt{write-config} read back
as it describes itself; three mutations detected. In the VM, a collection
driven by a \fichier{wac.yml}, without any option, is checked against the file
recorded in \fichier{investigation.json} (\texttt{check-json.py
--configuration}).

\section{Third-party chains: \texttt{third\_party\_chain\_test}}

Each rule of the check of third-party chains, on chains made by openssl to
break it (\fichier{third\_party\_chain\_fixtures.sh}): a sound chain accepted;
the root absent, excluded from code signing, distrusted; an authority's key
disallowed by Microsoft; an authority revoked according to the CCADB; an
issuer without \texttt{cA}; a signer for the web only; an authority with the
right name and another key; an ECDSA authority (not verified). Then the time
a chain is judged at, with a real RFC~3161 token made by \texttt{openssl ts}
and the collection's time moved to 2040: a signer valid 2020 to 2035 accepted
without time stamp today, refused in 2040, accepted in 2040 when stamped
today; a token over another signature, or altered, refused; a time stamping
authority under a root not trusted for it, refused; a root distrusted after
2030 accepted for a signature stamped before, refused otherwise. Last,
revocation, with real CRLs made by \texttt{openssl ca}: a sound chain accepted
with the authority's and the root's CRLs; without them, with the authority's
missing or forged (right issuer name, another key), refused as not
verifiable; a CRL found through the CCADB when the certificate's address is
missing, the certificate's own prevailing; a signer revoked for a compromise
refused even when stamped before; one revoked as superseded after the time
stamp accepted when stamped, refused otherwise. Each rule disabled in turn
fails the test (twenty-two mutations). The time stamps of the
real installers agree to the second with an independent reading
(\texttt{openssl ts} for RFC~3161 tokens, DER for counter-signatures). Real signatures are
cross-checked by \texttt{authenticode\_test --trust}, against \texttt{openssl
verify} with the set's \fichier{roots.pem} (installers of Git, Notepad++,
PuTTY, Python, VLC; the whole set and the set without the Sectigo roots); in
the VM, no chain \wac{} trusts is untrusted by
\texttt{Get-AuthenticodeSignature}.

\begin{lstlisting}
./third_party_chain_fixtures.sh /tmp/chain
/tmp/third_party_chain_test /tmp/chain
\end{lstlisting}

\section{Vulnerable drivers: \texttt{version\_info\_test}, \texttt{vulnerable\_drivers\_test}}

\texttt{version\_info\_test} checks the reading of the version resource
(original name, internal name, product, description, version) against
Windows' (PowerShell's \texttt{VersionInfo} in the VM) on 58 binaries:
versions and texts identical, except the texts a localised binary takes from
its \fichier{.mui} satellite — which the blocklist does not read —, then
truncated and damaged resources (11,852 checks).
\texttt{vulnerable\_drivers\_test} takes each rule alone: an Authenticode or
file fingerprint, a denied authority in the chain, narrowed to the signer, the
WHQL manufacturer, the original name, the product, the version bounds; a
binary without version resource (the file rule does not apply, as with
Windows) or not read (it applies) — nine mutations all detected. On real
LOLDrivers samples, \wac{}'s Authenticode digest is their ``Authentihash'' (4
out of 4), and two vulnerable drivers whose chain holds are refused; in the
VM, a vulnerable driver placed there (never loaded) must be recognised.

\begin{lstlisting}
g++ -std=c++17 -g -fsanitize=address,undefined -I. pe_resource.cpp \
    version_info.cpp version_info_test.cpp -o /tmp/version_info_test
/tmp/version_info_test <reference> <folder>
\end{lstlisting}

\section{Trust lists: \texttt{inflate\_test}, \texttt{cab\_test}, \texttt{zip\_test}, \texttt{csv\_reader\_test}, \texttt{trust\_list\_test}}

\opt{update-trust} decompresses and reads what Microsoft publishes; an error
would give a valid and wrong output. Each link therefore has its independent
reference, and is built on Linux with AddressSanitizer and UBSan:

\begin{itemize}
  \item \texttt{inflate\_test}: DEFLATE, checked against zlib on 49 streams,
    then streams truncated and damaged at random (fixed seed);
  \item \texttt{cab\_test}: MSZIP cabinets, extraction identical byte for byte
    to \texttt{cabextract}'s on the real \fichier{authrootstl.cab} and
    \fichier{disallowedcertstl.cab} (1,134 checks);
  \item \texttt{zip\_test}: ZIP archives, extraction identical to Python's
    \texttt{zipfile} on Microsoft's real vulnerable driver list and two
    synthetic archives (stored, compressed, comment, folders, empty file); a
    damaged archive never returns an altered content, refused by its CRC-32
    (1,717 checks);
  \item \texttt{csv\_reader\_test}: the CCADB CSV report, read field by
    field as Python's \texttt{csv} module reads it (10,284 records, SHA-256
    of every field), then doubled quotes, line breaks in a field, CRLF, a
    quote left open, a short record;
  \item \texttt{trust\_list\_test}: the counts of \fichier{authroot.stl} (roots
    for code signing, for time stamping, and for code signing not distrusted),
    computed apart by a DER reading in Python; the certificates Chromium blocks
    (\fichier{net/data/ssl/blocklist}), recognised by
    \fichier{disallowedcert.stl}; a list with one byte changed, refused; both
    lists read by eight threads at once, the chain cache empty, as the
    analysis threads of \opt{collect} \opt{binary} do. Built again with
    \texttt{-fsanitize=thread} (run through \texttt{setarch -R}, which
    ThreadSanitizer needs on recent kernels), it reported eleven data races on
    that cache before its lock.
\end{itemize}

\begin{lstlisting}
g++ -std=c++17 -g -fsanitize=address,undefined -I. authenticode.cpp rsa.cpp \
    sha.cpp quickdigest5.cpp trust_list_test.cpp -o /tmp/trust_list_test
/tmp/trust_list_test authroot.stl 313 335 145 disallowedcert.stl blocklist 36
\end{lstlisting}

Reference (lists of 25~August~2026 and 5~September~2025): 80 checks, no
failure. This test established what \fichier{disallowedcert.stl} holds, which
Microsoft does not document: its 16-byte entries are the MD5 of the
\textbf{public key} (not of the signed part, as the property's name suggests),
its 48-byte entries the SHA-384 of the signed part.

The end-to-end result is checked in the VM: the 562 roots \opt{update-trust}
writes are identical, byte for byte, to those \texttt{certutil
-generateSSTFromWU} produces. The vulnerable driver lists are cross-checked by
a Python reading of the downloaded XML and JSON: the same 1,086 Microsoft
fingerprints, 217 denied signers, 8,151 LOLDrivers fingerprints. The CRLs
are checked with openssl: each one reads, and its issuer is the authority the
CCADB associates with it (387 out of 387).

% ============================================================ 5
\chapter{The robustness harnesses}\label{ch:robustesse}

\section{The principle}

A parser reads at offsets the file itself declares. On a truncated or forged
file, a read past the buffer does not crash on its own: it reads the
neighbouring memory and publishes it, in valid JSON. The harnesses make it
crash: each input is copied against a forbidden memory page
(\texttt{PAGE\_NOACCESS}), once \emph{ending} against it (reads past the end),
once \emph{starting} right after it (reads before the start, for instance at a
negative offset taken from the data). A single byte read outside the input
raises an access violation.

For each file given, the harness parses:
\begin{itemize}
  \item the whole file;
  \item its truncations (every length for small files);
  \item copies with bytes overwritten at random, with a fixed seed: a failure
    reproduces identically.
\end{itemize}

\section{Shortcuts: \texttt{lnk\_test}}

\begin{lstlisting}
cd build-windows/tests
wine lnk_test.exe a.lnk b.lnk ...
\end{lstlisting}

Where to find files: the \fichier{Recent} folder and the Start menu of the VM,
or those of a collection (\fichier{Users\textbackslash
<name>\textbackslash AppData\textbackslash Roaming\textbackslash
Microsoft\textbackslash Windows\textbackslash Recent}). Accented or non-Latin
names travel badly through the \texttt{wine} command line: copy the files under
numbered names.

\section{Jump lists and Prefetch: \texttt{parsers\_test}}\label{sec:parsers}

\begin{lstlisting}
parsers_test.exe jumplist-auto   <file.automaticDestinations-ms> ...
parsers_test.exe jumplist-custom <file.customDestinations-ms> ...
parsers_test.exe prefetch        <file.pf> ...
parsers_test.exe hive            <hive> ...
\end{lstlisting}

The \texttt{hive} kind tries the hive reader: each input is opened from memory,
against the forbidden page, and the whole tree is walked (subkeys, values,
\texttt{ORQueryInfoKey}). A subkey pointing back at one of its ancestors is
refused by the reader: without that, a recursive walk would never end.

Jump lists are tested under \texttt{wine}, like \texttt{lnk\_test}. Prefetch files
are tested \textbf{in the VM}: they are compressed (XPRESS Huffman), and only the
\fichier{ntdll} library of Windows decompresses them. The harness decompresses
each file before trying it: it is \wac{}'s parser it aims at, not Windows'
decompressor.

For the 307 Prefetch files of a real machine (over a million inputs), count about
two hours in the VM: run the test as a detached task, in batches, and follow its
log. Example PowerShell script, put into the VM with \texttt{qga.py write}:

\begin{lstlisting}
$files = Get-ChildItem C:\wactest\pt\pf\*.pf | ForEach-Object { $_.FullName }
$log = 'C:\wactest\pt\pf_results.txt'
for ($i = 0; $i -lt $files.Count; $i += 20) {
  $batch = $files[$i..([Math]::Min($i + 19, $files.Count - 1))]
  Add-Content $log (& C:\wactest\pt\parsers_test.exe prefetch @batch 2>&1)
  Add-Content $log ("batch " + $i + " exit=" + $LASTEXITCODE)
}
Add-Content $log 'DONE'
\end{lstlisting}

\section{Reference}

Results obtained on a real machine, each input parsed on both sides of the
forbidden page:

\begin{longtable}{@{}L{6cm}L{3.2cm}L{5.4cm}@{}}
\toprule
\textbf{Files} & \textbf{Inputs} & \textbf{Where} \\
\midrule
\endhead
220 shortcuts & 546,214 & \texttt{wine} \\
48 automatic jump lists & 175,585 & \texttt{wine} \\
23 custom jump lists & 79,036 & \texttt{wine} \\
307 Prefetch files & 1,246,491 & VM, about 2~h \\
4 registry hives & 16,207 & \texttt{wine} \\
\bottomrule
\end{longtable}

% ============================================================ 6
\chapter{Cross-checking on a real collection}

A freshly installed VM produces no fragmentation, no old jump lists, no exotic
file names. Several defects only showed on a real collection, confronted with an
independent source:

\begin{itemize}
  \item \textbf{jump lists / \texttt{olefile}}: the streams and the DestList,
    read again by an independent OLE reader, must give exactly the same entries
    (numbers, dates, paths). That is what showed the streams are named ``1'' to
    ``f'', not ``01'' to ``0f'';
  \item \textbf{shortcut / DestList}: in a jump list, the target of a shortcut
    and the path of its DestList entry are two records of the same file. Their
    disagreement revealed a wrong code page (``d\textbackslash x92orientation'')
    and an ignored Unicode path;
  \item \textbf{events / \texttt{python-evtx}}: record-by-record comparison
    (\texttt{evtx\_test --dump});
  \item \textbf{signatures / \texttt{Get-AuthenticodeSignature}}
    (\texttt{authenticode\_test});
  \item \textbf{old collection / new one}: decode the same files with the
    current code and compare, after normalisation, with the output of a
    previous version.
\end{itemize}

\begin{attention}
A real collection is processed locally. If a hive must be confronted with
Microsoft's reader (\texttt{offline\_registry\_test}), it may be copied into
the test VM, then deleted as soon as the test is over. No copy stays in the temporary working folder, and none of it
enters the repository.
\end{attention}

% ============================================================ 7
\chapter{Reading a failure}

\section{Exit codes}

\begin{longtable}{@{}L{4.4cm}L{10.6cm}@{}}
\toprule
\textbf{Code} & \textbf{Meaning} \\
\midrule
\endhead
0 & everything passed \\
1 & a check failed, or an input file is unreadable \\
2 & wrong usage (arguments) \\
3 & \texttt{parsers\_test}: a parse did not end within 60~s (watchdog) \\
\texttt{0xC0000005} (-1073741819) & access violation: a read outside the buffer \\
\bottomrule
\end{longtable}

\section{Locating a read outside the buffer}

Under \texttt{wine}, a fault prints the call stack, with addresses relative to
the program (\texttt{lnk\_test+0x485bb}). On Windows, \texttt{parsers\_test}
itself prints the input at fault and the faulting address:

\begin{lstlisting}
FAILED  exception 0xC0000005 at exe+0x78b70, address read 0x159deb4200c:
        C:\wactest\pt\pf\CMD.EXE-0BD30981.pf truncated at 192 (guard after)
\end{lstlisting}

The address is translated into a function with the program's symbol table (base
address \texttt{0x140000000}):

\begin{lstlisting}
x86_64-w64-mingw32-nm -C parsers_test.exe | grep -i " t " | sort > syms.txt
# look for the symbol whose address is just below 0x140000000 + 0x78b70
x86_64-w64-mingw32-objdump -d --start-address=0x140078b40 \
    --stop-address=0x140078b80 parsers_test.exe
\end{lstlisting}

The disassembled instruction gives the offset read (here \texttt{0xcc(\%r14)});
it is then matched, in the source code, with the field read at that position.
Then replay the exact input (the truncation or corruption shown) to check the
fix, and the whole series after it.

\section{An endless loop}

If \texttt{parsers\_test} stops on ``no end after 60~s'', the input named keeps a
parser running for ever — typically a cyclic sector chain in an OLE container.
It is a defect just like a read outside the buffer: a forged file would be
enough to block a collection.

\end{document}
